\chapter{Основные причины возникновения рассинхронизации}

Рассинхронизация состояния - это расхождение между данными, хранящимися на сервере, и состоянием,
которое отображается на клиентах. Возникновение рассинхронизации обусловлено как особенностями сетевого
окружения, так и архитектурными проблемами внутри игрового движка.

\section{Сетевые задержки}

Одной из наиболее распространённых причин рассинхронизации является задержка передачи данных\cite{pantel2002lag}.
Игровой клиент всегда получает состояние с определённой задержкой, зависящей от качества соединения.
Если задержка непостоянна и резко меняется, клиент может получить обновления позже или раньше
ожидаемого момента, что вызывает временные расхождения в позициях объектов.

Ситуация ухудшается в условиях плохих сетей, где вариативность задержек возрастает.

\section{Потеря, дублирование и нарушение порядка доставки пакетов}

При использовании ненадёжных протоколов часть пакетов может теряться или приходить в другом порядке.
Это приводит к следующим ошибкам:

\begin{itemize}
    \item объект остаётся на устаревшей позиции;
    \item проигрываются неактуальные события;
    \item происходит откат состояния назад.
\end{itemize}

Хотя игровые движки реализуют механизмы частичной надёжности, абсолютной защиты от подобных проблем нет.

\section{Недетерминизм симуляции}

Некоторые игры опираются на локальный расчёт физики или игровой логики на стороне клиента. Любое незначительное
отличие в вычислениях может привести к расхождению состояния. Основные источники недетерминизма:

\begin{itemize}
    \item использование чисел с плавающей точкой, дающих разные результаты на разных процессорах;
    \item различный порядок обновления объектов на разных клиентах;
    \item вариации частоты кадров, влияющие на расчёт физики;
\end{itemize}

Такие проблемы особенно заметны в RTS и симуляторах.

\section{Различия в производительности клиентских устройств}

Игровые клиенты могут работать с разной скоростью в зависимости от мощности процессора или видеокарты.
Если частота обновления симуляции отличается:

\begin{itemize}
    \item одни клиенты рассчитывают больше физических шагов, чем другие;
    \item предсказания состояния расходятся с серверной логикой.
\end{itemize}

В играх, где симуляция сильно зависит от стабильной частоты кадров, такие различия являются источником ошибок.

\section{Неправильная или неполная сериализация данных}

При передаче состояния объекты сериализуются.
Ошибки сериализации приводят к тому, что данные не совпадают на клиенте и сервере. Это может быть вызвано:

\begin{itemize}
    \item пропуском некоторых параметров объекта;
    \item отправкой данных в неправильном формате;
    \item несогласованностью типов при чтении и записи.
\end{itemize}

\section{Несогласованный порядок выполнения событий}

Если события обрабатываются не в том же порядке, что на сервере, результаты симуляции расходятся.
Это происходит из-за:

\begin{itemize}
    \item локальных оптимизаций клиента;
    \item попадания пакетов с разными временными метками;
    \item параллельного выполнения задач;
    \item применения предсказания до получения подтверждения от сервера.
\end{itemize}

Несогласованная обработка событий особенно критична для игр с динамичными боевыми механиками.

\section{Ошибки в предсказании и последующей коррекции состояния}

Client-side prediction значительно повышает отзывчивость, но несёт риск ошибок.
Предсказание может расходиться с реальными данными сервера, что приводит к:

\begin{itemize}
    \item резким скачкам при корректировке\cite{gamenetpatterns};
    \item неверным расчётам следующего состояния;
    \item накоплению ошибки при плохой связи.
\end{itemize}

Такие проблемы заметны при нестабильной задержки сети и большом количестве корректировок.

